% !PRINTED 96
\weberpage{96}
%<<CONT

\[
\begin{matrix}
 A_1^{(1)}&A_2^{(1)}&\cdots&A_n^{(1)}\\
 A_1^{(2)}&A_2^{(2)}&\cdots&A_n^{(2)}\\
 \cdots&\cdots&&\cdots\\
 A_1^{(n)}&A_2^{(n)}&\cdots&A_n^{(n)}
\end{matrix}\tag{2}
\]
das System der Unterdeterminanten. Bilden wir aus (2) die Determinante, die wir mit $\Delta$ bezeichnen wollen, so können wir auf das Product $A\Delta$ die Multiplicationsregel anwenden. Dies giebt aber nach §. 22, (3) und (7):
\[
 A\Delta=\begin{vmatrix}
 A&0&\cdots&0\\
 0&A&\cdots&0\\
 \cdots&\cdots&&\cdots\\
 0&0&\cdots&A
\end{vmatrix}=A^n,
\]
und daraus durch Division mit $A$:
\begin{equation}
 \Delta=A^{n-1}.\tag{3}
\end{equation}
Es ist also $\Delta$ die $(n-1)$te Potenz von $A$. Bei dieser Ableitung ist allerdings zunächst vorausgesetzt, dass $A$ von Null verschieden sei. Da aber (3) in Bezug auf die Elemente $a_i^{(h)}$ eine Identität ist, d. h. auch dann gilt, wenn diese Grössen unabhängige Variable sind, so folgt, dass auch noch in diesem Ausnahmefall die Formel (3) gilt, d. h. dass, wenn $A$ verschwindet, auch $\Delta$ verschwindet.

Dies Ergebniss ist ein specieller Fall eines allgemeineren Satzes, nach dem jede beliebige Determinante der Matrix (2) gebildet werden kann. Betrachten wir die $\nu$-reihige Unterdeterminante
\[
\Delta_\nu=\begin{vmatrix}
 A_1^{(1)}&A_2^{(1)}&\cdots&A_\nu^{(1)}\\
 A_1^{(2)}&A_2^{(2)}&\cdots&A_\nu^{(2)}\\
 \cdots&\cdots&&\cdots\\
 A_1^{(\nu)}&A_2^{(\nu)}&\cdots&A_\nu^{(\nu)}
\end{vmatrix},\tag{4}
\]
aus der man durch Permutation der oberen und unteren Indices alle anderen $\nu$-reihigen Unterdeterminanten ableiten kann, so kann man die Multiplicationsregel anwenden, indem man $\Delta_\nu$ nach der Schlussbemerkung des letzten Paragraphen in eine $n$-reihige Determinante verwandelt.
